\section{Data Structure}
\subsection{PBDS}
\begin{minted}{c++}
#include <ext/pb_ds/assoc_container.hpp>
#include <ext/pb_ds/tree_policy.hpp>
using namespace __gnu_pbds;
template <class T, class Cmp = less<T>>
using OS = tree<T, null_type, Cmp, rb_tree_tag,
                tree_order_statistics_node_update>;

// Multiset: strictly use pair<T, int> to avoid UB!
// Insert: s.insert({val, id++})
// Erase exactly x: s.erase(s.lower_bound({x, 0}))

// find_by_order(k) -> iterator to k-th element
// order_of_key({x, 0}) -> count of elements < x

// Lower Bound of x (first element >= x):
// int idx = s.order_of_key({x, 0}); // 0-indexed
// if (idx == s.size()) no_lower_bound;
// else lb = s.find_by_order(idx)->first;

// Max element < x: s.find_by_order(idx - 1)->first
// Upper Bound of x: s.order_of_key({x + 1, 0})

// --- Custom Comparator Example ---
struct cmp {
  bool operator()(int a, int b) { return a > b; } // must be strict
};
using custom_set = OS<int, cmp>;
\end{minted}
\subsection{GP Hash Table}
\textbf{Usage:} \texttt{gp\_hash\_table<int, int, custom\_hash> mp;} (Drop-in replacement for \texttt{unordered\_map} to prevent anti-hash tests).
\begin{minted}{c++}
struct custom_hash {
  static uint64_t splitmix64(uint64_t x) {
    x += 0x9e3779b97f4a7c15;
    x = (x ^ (x >> 30)) * 0xbf58476d1ce4e5b9;
    x = (x ^ (x >> 27)) * 0x94d049bb133111eb;
    return x ^ (x >> 31);
  }
  size_t operator()(uint64_t x) const {
    static const uint64_t FIXED_RANDOM =
    chrono::steady_clock::now().time_since_epoch()
    .count();
    return splitmix64(x + FIXED_RANDOM);
  }
  size_t operator()(
    const pair<ll, ll>& p) const noexcept {
    uint64_t a =
    splitmix64(p.first + 0x9e3779b97f4a7c15ULL);
    uint64_t b = splitmix64(p.second +
    0x9e3779b97f4a7c15ULL + 0xC6A4A7935BD1E995ULL);
    return (size_t)(a ^ (b << 1));
  }
};
gp_hash_table<pair<ll,ll>, ll, custom_hash> mp;
\end{minted}
\subsection{Rope (GNU)}
\texttt{gnu\_cxx::rope} = built-in treap (split/merge/insert/erase $O(\log N)$). \texttt{crope} for char, \texttt{rope<int>} for ints. Include:
\begin{minted}{c++}
#include <ext/rope>
using namespace __gnu_cxx;
crope r; // crope for char, rope<int> for ints
r.push_back('a');            // append
r.insert(0,"hello ");        // insert string
r.erase(0,6);                // erase 6 chars from pos 0
r.substr(0,5);               // extract substring [0,5)
char c = r[3];               // access O(log N)
cout << r << "\n";           // print whole rope
string s="test"; crope r2(s.begin(),s.end());
\end{minted}
\subsection{DSU}
KACTL-style single-vector DSU: \texttt{e[x]} is negative for roots (stores \texttt{-size}) and non-negative for parents. $66\%$ memory savings vs 3-vector approach. \texttt{size(x)} returns component size. \texttt{cnt} tracks number of components.
\begin{minted}{c++}
struct dsu {
  vector<int> e;  // negative = size (root), non-negative = parent
  int cnt;
  dsu(int n) : e(n, -1), cnt(n) {}
  int find(int x) {
    return e[x] < 0 ? x : e[x] = find(e[x]);
  }
  bool unite(int x, int y) {
    x = find(x), y = find(y);
    if (x == y) return false;
    if (e[x] > e[y]) swap(x, y);  // more negative = larger size
    e[x] += e[y]; e[y] = x; --cnt;
    return true;
  }
  int size(int x) { return -e[find(x)]; }
};
\end{minted}
\subsection{Sparse Table}
Range query (max/min/GCD). Idempotent --- overlapping intervals OK. Build $O(N \log N)$, query $O(1)$. No updates. Use log-major layout for cache efficiency.
\begin{minted}{c++}
struct Sparse {
  vector<vector<int>> t;
  Sparse(const vector<int>& a) {
    int n = a.size(), lg = n ? 32 - __builtin_clz(n) : 0;
    t.assign(lg, vector<int>(n));
    t[0] = a;
    for (int j = 1; j < lg; j++)
      for (int i = 0; i + (1 << j) <= n; i++)
        t[j][i] = max(t[j-1][i], t[j-1][i+(1<<(j-1))]);
  }
  int qry(int l, int r) const {
    int j = 31 - __builtin_clz(r-l+1);
    return max(t[j][l], t[j][r-(1<<j)+1]);
  }
};
// For min: s/max/min/   For gcd: s/max/gcd/
\end{minted}
\subsection{Sparse Table 2D}
Build $O(NM \log N \log M)$, query $O(1)$. \textbf{Memory intensive:} $\sim 600$ MB for $2048^2$ grid (11$\times$11 matrices). Use for small grids only. Prefer 2D Fenwick for larger grids.
\begin{minted}{c++}
struct Sparse2D {
  int n, m;
  vector<vector<int>> a;
  vector<int> lg;
  vector<vector<int>> st[11][11];  // st[logA][logB] is a matrix; supports up to 2047
  Sparse2D(int n, int m) : n(n), m(m) {
    a.assign(n, vector<int>(m));
    lg.resize(max(n, m) + 1);
    for (int i = 2; i <= max(n, m); i++) lg[i] = lg[i>>1] + 1;
  }
  void set_val(int i, int j, int v) { a[i][j] = v; }
  void build() {
    auto &s = st[0][0];
    s.assign(n, vector<int>(m));
    for (int i = 0; i < n; i++)
      for (int j = 0; j < m; j++)
        s[i][j] = a[i][j];
    int la = lg[n], lb = lg[m];
    for (int ka = 0; ka <= la; ka++) {
      for (int kb = 0; kb <= lb; kb++) {
        if (ka + kb == 0) continue;
        int pwa = 1 << ka, pwb = 1 << kb;
        auto &cur = st[ka][kb];
        cur.assign(n - pwa + 1, vector<int>(m - pwb + 1));
        for (int i = 0; i + pwa <= n; i++)
          for (int j = 0; j + pwb <= m; j++) {
            if (!ka)
              cur[i][j] = max(st[ka][kb-1][i][j],
                st[ka][kb-1][i][j+pwb/2]);
            else
              cur[i][j] = max(st[ka-1][kb][i][j],
                st[ka-1][kb][i+pwa/2][j]);
          }
      }
    }
  }
  int qry(int x1, int y1, int x2, int y2) const {
    int a = lg[x2-x1+1], b = lg[y2-y1+1];
    auto &s = st[a][b];
    int m1 = max(s[x1][y1], s[x2-(1<<a)+1][y1]);
    int m2 = max(s[x1][y2-(1<<b)+1],
                 s[x2-(1<<a)+1][y2-(1<<b)+1]);
    return max(m1, m2);
  }
};
\end{minted}
\subsection{Fenwick Tree}
Point update, prefix/range sum. $O(\log N)$ per op. 1-indexed internally. Usage: \texttt{BIT<ll> fenw(n); fenw.add(i, val); fenw.sum(l, r);}
\begin{minted}{c++}
template<typename T>
struct BIT {
  int n;
  vector<T> bit;
  BIT(int n = 0) { init(n); }
  void init(int n_) { n = n_; bit.assign(n + 1, 0); }
  void add(int i, T v) {
    // walk up: i += i & -i moves to next responsible ancestor
    for (; i <= n; i += i & -i) bit[i] += v; }
  T sum(int i) {
    T r = 0;
    // walk down: i -= i & -i strips lowest set bit
    for (; i > 0; i -= i & -i) r += bit[i];
    return r;
  }
  T sum(int l, int r) { return sum(r) - sum(l - 1); }
  // find smallest i with prefix_sum(i) >= k (values must be non-negative)
  int find_kth(T k) {
    int i = 0;
    for (int j = 1 << (31 - __builtin_clz(n)); j; j >>= 1)
      if (i + j <= n && bit[i + j] < k) k -= bit[i += j];
    return i + 1;
  }
};
\end{minted}
Usage: \texttt{BIT<ll> fenw(5); fenw.find\_kth(k)} returns 1-indexed position.
\subsection{Fenwick 2D}
\begin{minted}{c++}
struct BIT2D {
  int n, m;
  vector<vector<ll>> bit;
  BIT2D(int n, int m): n(n), m(m) {
    bit.assign(n + 1, vector<ll>(m + 1, 0));
  }
  void add(int x, int y, ll val) {
    for (int i = x; i <= n; i += i & -i)
      for (int j = y; j <= m; j += j & -j)
        bit[i][j] += val;
  }
  ll sum(int x, int y) {
    ll res = 0;
    for (int i = x; i > 0; i -= i & -i)
      for (int j = y; j > 0; j -= j & -j)
        res += bit[i][j];
    return res;
  }
  ll query(int x1, int y1, int x2, int y2) {
    return sum(x2, y2) - sum(x1 - 1, y2) -
    sum(x2, y1 - 1) + sum(x1 - 1, y1 - 1);
  }
};
BIT2D fenw2(4, 4);
fenw2.add(1, 1, 5);
fenw2.add(3, 2, 7);
cout << fenw2.query(1, 1, 3, 3) << "\n";
\end{minted}
\subsection{BIT Range Update \texorpdfstring{\&}{&} Range Query}
Uses 2 BITs trick. Add val to $[l,r]$, query sum of $[l,r]$.
Derive from prefix sum identity: $psum(i) = i \cdot bit_1(i) - bit_2(i)$.
\begin{minted}{c++}
struct BITRange {
  int n;
  vector<ll> bit1, bit2;
  BITRange(int n) : n(n), bit1(n + 2), bit2(n + 2) {}
  void add(vector<ll>& bit, int i, ll v) {
    for (; i <= n; i += i & -i) bit[i] += v;
  }
  void upd(int l, int r, ll v) {
    add(bit1, l, v);
    add(bit1, r + 1, -v);
    add(bit2, l, v * (l - 1));
    add(bit2, r + 1, -(r * v));
  }
  ll sum(vector<ll>& bit, int i) {
    ll r = 0;
    for (; i > 0; i -= i & -i) r += bit[i];
    return r;
  }
  ll qry(int l, int r) {
    auto pref = [&](int i) { return sum(bit1, i) * i - sum(bit2, i); };
    return pref(r) - pref(l - 1);
  }
};
\end{minted}
\texttt{BITRange bit(n);} \hfill \texttt{bit.upd(l,r,v);} \hfill \texttt{bit.qry(l,r);}
\subsection{Lazy Segment Tree}
Range add/set, range sum/max/min. $O(\log N)$ per op. Uses \texttt{NO\_SET} sentinel to eliminate the \texttt{hz} flag vector -- saves $\sim 2.4\%$ memory. Requires \texttt{using ll = long long;} in setup.
\begin{minted}{c++}
struct LSeg {
  int n;
  vector<ll> sm, mx, mn, ad, st;
  const ll INF = 1e18, NO_SET = 9e18;
  LSeg(int n): n(n), sm(4*n), mx(4*n), mn(4*n),
    ad(4*n), st(4*n, NO_SET) {}
  void pull(int i) {
    sm[i]=sm[i*2]+sm[i*2+1];
    mx[i]=max(mx[i*2],mx[i*2+1]);
    mn[i]=min(mn[i*2],mn[i*2+1]);
  }
  void pset(int i,int s,int e,ll v) {
    st[i]=v; ad[i]=0;
    sm[i]=v*(e-s+1); mx[i]=mn[i]=v;
  }
  void padd(int i,int s,int e,ll v) {
    if(st[i]!=NO_SET)st[i]+=v; else ad[i]+=v;
    sm[i]+=v*(e-s+1); mx[i]+=v; mn[i]+=v;
  }
  void push(int i,int s,int e) {
    if(s==e)return;
    int m=(s+e)/2;
    if(st[i]!=NO_SET){pset(i*2,s,m,st[i]);
      pset(i*2+1,m+1,e,st[i]);st[i]=NO_SET;}
    if(ad[i]){padd(i*2,s,m,ad[i]);
      padd(i*2+1,m+1,e,ad[i]);ad[i]=0;}
  }
  void upd(int l,int r,ll v,int t,
    int i=1,int s=0,int e=-1) {
    if(e<0)e=n-1;
    if(r<s||e<l)return;
    if(l<=s&&e<=r)
      {t?pset(i,s,e,v):padd(i,s,e,v);return;}
    push(i,s,e);int m=(s+e)/2;
    upd(l,r,v,t,i*2,s,m);
    upd(l,r,v,t,i*2+1,m+1,e);
    pull(i);
  }
  array<ll,3> qry(int l,int r,
    int i=1,int s=0,int e=-1) {
    if(e<0)e=n-1;
    if(r<s||e<l)return {0,-INF,INF};
    if(l<=s&&e<=r)return {sm[i],mx[i],mn[i]};
    push(i,s,e);int m=(s+e)/2;
    auto a=qry(l,r,i*2,s,m),
         b=qry(l,r,i*2+1,m+1,e);
    return {a[0]+b[0],max(a[1],b[1]),
            min(a[2],b[2])};
  }
  void set(int l,int r,ll v){upd(l,r,v,1);}
  void add(int l,int r,ll v){upd(l,r,v,0);}
  // auto [s,mx,mn]=seg.qry(l,r);
  // seg.set(l,r,v); seg.add(l,r,v);
};
\end{minted}
\subsection{Persistent Segment Tree}
Persistent segment tree supporting k-th smallest, count $\le x$, and range sum (distinct elements). All operations $O(\log N)$. Each version = new root.

\textbf{Usage 1 -- k-th smallest / count $\le x$ (value tree):} \\
\texttt{root[0] = build(1, m);}\quad($m$ = \#unique values) \\
\texttt{for i=1..n: root[i] = upd(root[i-1], 1, m, comp[a[i]], +1)} \\
\texttt{kth(root[l-1], root[r], 1, m, k)} \hfill $k$-th smallest in $[l,r]$ \\
\texttt{cnt(root[l-1], root[r], 1, m, x)} \hfill $\# \le x$ in $[l,r]$ \\
\texttt{cnt(root[l-1], root[r], y) - cnt(x-1)} \hfill $\# \in [x,y]$ \\

\textbf{Usage 2 -- distinct count (position tree):} \\
\texttt{root[0] = build(1, n);} \\
\texttt{for i=1..n:} \\
\qquad\texttt{r = root[i-1];} \\
\qquad\texttt{if (last[a[i]]) r = upd(r, 1, n, last[a[i]], -1);} \\
\qquad\texttt{root[i] = upd(r, 1, n, i, +1);} \\
\qquad\texttt{last[a[i]] = i;} \\
\texttt{sum(root[0], root[r], 1, n, l, r)} \hfill distinct count in $[l,r]$

\begin{minted}{c++}
struct PST {
  static const int N = 2'00'000;
  static const int LG = 18;
  static const int M = 1 + 4 * N + N * LG;
  int lc[M], rc[M], val[M];
  int sz = 1;

  int build(int b, int e) {
    int cur = sz++;
    if (b == e) return cur;
    int m = (b + e) >> 1;
    lc[cur] = build(b, m);
    rc[cur] = build(m + 1, e);
    return cur;
  }
  int upd(int pre, int b, int e, int i, int v) {
    int cur = sz++;
    lc[cur] = lc[pre]; rc[cur] = rc[pre];
    if (b == e) { val[cur] = val[pre] + v; return cur; }
    int m = (b + e) >> 1;
    if (i <= m) lc[cur] = upd(lc[pre], b, m, i, v);
    else        rc[cur] = upd(rc[pre], m + 1, e, i, v);
    val[cur] = val[lc[cur]] + val[rc[cur]];
    return cur;
  }
  int kth(int pre, int cur, int b, int e, int k) {
    while (b < e) {
      int m = (b + e) >> 1;
      int left_cnt = val[lc[cur]] - val[lc[pre]];
      if (k <= left_cnt) { cur = lc[cur]; pre = lc[pre]; e = m; }
      else               { cur = rc[cur]; pre = rc[pre]; b = m + 1; k -= left_cnt; }
    }
    return b;
  }
  int cnt(int pre, int cur, int b, int e, int x) {
    if (x < b) return 0;
    int res = 0;
    while (b < e) {
      int m = (b + e) >> 1;
      if (x <= m) { cur = lc[cur]; pre = lc[pre]; e = m; }
      else { res += val[lc[cur]] - val[lc[pre]]; cur = rc[cur]; pre = rc[pre]; b = m + 1; }
    }
    res += val[cur] - val[pre];
    return res;
  }
  int sum(int pre, int cur, int b, int e, int l, int r) {
    if (l > e || r < b) return 0;
    if (l <= b && e <= r) return val[cur] - val[pre];
    int m = (b + e) >> 1;
    return sum(lc[pre], lc[cur], b, m, l, r) + sum(rc[pre], rc[cur], m + 1, e, l, r);
  }
};
\end{minted}
\subsection{Merge Sort Tree}
Segment tree where each node stores a sorted list of its range. $qry(l,r,k)$ = number of elements $\le k$ in $[l,r]$. Build $O(N \log N)$, query $O(\log^2 N)$, space $O(N \log N)$.
\begin{minted}{c++}
struct MergeSortTree {
  int n;
  vector<vector<int>> tree;
  MergeSortTree(const vector<int>& arr) {
    n = arr.size();
    tree.resize(4 * n);
    build(1, 0, n - 1, arr);
  }
  void build(int node, int l, int r,
             const vector<int>& arr) {
    if (l == r) {
      tree[node].push_back(arr[l]);
      return;
    }
    int mid = (l + r) / 2;
    build(2 * node, l, mid, arr);
    build(2 * node + 1, mid + 1, r, arr);
    merge(tree[2 * node].begin(), tree[2 * node].end(),
          tree[2 * node + 1].begin(),
          tree[2 * node + 1].end(),
          back_inserter(tree[node]));
  }
  int query(int node, int l, int r, int ql, int qr, int k) {
    if (l > qr || r < ql) return 0;
    if (l >= ql && r <= qr) {
      // binary search in sorted node list: count <= k
      return upper_bound(tree[node].begin(),
      tree[node].end(), k) - tree[node].begin();
    }
    int mid = (l + r) / 2;
    return query(2 * node, l, mid, ql, qr, k) +
           query(2 * node + 1, mid + 1, r, ql, qr, k);
  }
};
\end{minted}
% PST: O(log N) with compression. MST: O(log^2 N), no compression needed.

\subsection{Wavelet Matrix}
Wavelet Matrix (compact iterative WT) using \texttt{vector<uint64\_t>} bit vectors instead of \texttt{vector<int>} prefix arrays --- $32\times$ memory savings per level. Build $O(N \log M)$, query $O(\log M)$. 0-indexed. \texttt{kth(l,r,k)} returns $k$-th smallest (0-based) in $[l,r]$. \texttt{count\_leq(l,r,x)} counts $\le x$. \texttt{count\_seg(l,r,x,y)} counts in $[x,y]$.
\begin{minted}{c++}
struct WT {
  int n, lg;
  vector<vector<uint64_t>> b;
  vector<vector<int>> p; // prefix sum over 64-bit blocks
  vector<int> z;         // number of 0s at each level
  WT(vector<int> v) : n(v.size()) {
    int mx = n ? *max_element(v.begin(), v.end()) : 0;
    lg = mx ? 32 - __builtin_clz(mx) : 0;
    if (!lg) { b.resize(1); p.resize(1); z.resize(1); return; }
    b.resize(lg); p.resize(lg); z.resize(lg);
    vector<int> cur = v, nxt(n);
    for (int d = lg - 1; d >= 0; --d) {
      int bc = (n + 63) / 64;
      b[d].assign(bc, 0);
      int p0 = 0, p1 = 0;
      for (int i = 0; i < n; ++i) {
        if (cur[i] >> d & 1) {
          b[d][i / 64] |= 1ULL << (i % 64);
          nxt[p1++] = cur[i];
        } else cur[p0++] = cur[i];
      }
      z[d] = p0;
      copy(nxt.begin(), nxt.begin() + p1, cur.begin() + p0);
      p[d].resize(bc + 1);
      for (int i = 0; i < bc; ++i)
        p[d][i + 1] = p[d][i] + __builtin_popcountll(b[d][i]);
    }
  }
  int rank1(int d, int i) const {
    int bc = (n + 63) / 64;
    if (i >= n) return p[d][bc];
    int r = i / 64;
    return p[d][r] + __builtin_popcountll(
      b[d][r] & ((1ULL << (i % 64)) - 1));
  }
  int rank0(int d, int i) const { return i - rank1(d, i); }
  int kth(int l, int r, int k) {
    int ans = 0;
    for (int d = lg - 1; d >= 0; --d) {
      int c0 = rank0(d, r) - rank0(d, l);
      if (k < c0) {
        l = rank0(d, l); r = rank0(d, r);
      } else {
        ans |= 1 << d; k -= c0;
        l = z[d] + rank1(d, l); r = z[d] + rank1(d, r);
      }
    }
    return ans;
  }
  int count_leq(int l, int r, int x) {
    if (x < 0) return 0;
    if (r <= l) return 0;
    if (x >= (1LL << lg)) return r - l;
    int ans = 0;
    for (int d = lg - 1; d >= 0; --d) {
      int c0 = rank0(d, r) - rank0(d, l);
      if (x >> d & 1) {
        ans += c0;
        l = z[d] + rank1(d, l); r = z[d] + rank1(d, r);
      } else {
        l = rank0(d, l); r = rank0(d, r);
      }
    }
    return ans + (r - l);
  }
  int count_seg(int l, int r, int x, int y) {
    if (x > y) return 0;
    return count_leq(l, r, y) - count_leq(l, r, x - 1);
  }
};
\end{minted}
\subsection{Trie (Binary)}
Binary Trie for integers up to $2^{B}-1$. $O(B)$ per op. Uses single \texttt{vector<Node>} pool for cache locality (all node data in one cache line vs two separate vectors).
\texttt{add(x, v)} -- insert ($v=1$) or erase ($v=-1$)\\
\texttt{count(x)} -- how many times $x$ inserted\\
\texttt{count\_xor\_less(x, k)} -- $\#\{y : (x \oplus y) < k\}$\\
\texttt{count\_xor\_greater(x, k)} -- $\#\{y : (x \oplus y) > k\}$\\
\texttt{get\_kth(k)} -- $k$-th smallest value (1-indexed)\\
\texttt{get\_xor\_kth(x, k)} -- $k$-th smallest $(x \oplus y)$ over $y$\\
\texttt{count\_special(x)} -- $\#\{y : y \le (x \oplus y) \le x\}$
\begin{minted}{c++}
struct Trie {
  static const int B = 31;
  struct Node { int nxt[2]; int cnt; };
  vector<Node> pool;
  Trie() { pool.push_back({0, 0, 0}); }
  void add(int x, int v = 1) {
    int u = 0; pool[u].cnt += v;
    for (int i = B - 1; i >= 0; --i) {
      int d = (x >> i) & 1;
      if (!pool[u].nxt[d]) {
        pool[u].nxt[d] = pool.size();
        pool.push_back({0, 0, 0});
      }
      u = pool[u].nxt[d]; pool[u].cnt += v;
    }
  }
  int count(int x) {
    int u = 0;
    for (int i = B - 1; i >= 0; --i) {
      int d = (x >> i) & 1;
      int nxt = pool[u].nxt[d];
      if (!nxt || !pool[nxt].cnt) return 0;
      u = nxt;
    }
    return pool[u].cnt;
  }
  int count_xor_less(int x, int k) {
    int u = 0, ans = 0;
    for (int i = B - 1; i >= 0; --i) {
      int d = (x >> i) & 1;
      if ((k >> i) & 1) {
        if (pool[u].nxt[d]) ans += pool[pool[u].nxt[d]].cnt;
        if (!pool[u].nxt[d ^ 1]) break;
        u = pool[u].nxt[d ^ 1];
      } else {
        if (!pool[u].nxt[d]) break;
        u = pool[u].nxt[d];
      }
    }
    return ans;
  }
  int count_xor_greater(int x, int k) {
    if (k >= (1LL << B) - 1) return 0;
    return pool[0].cnt - count_xor_less(x, k + 1); }
  int get_kth(int k) { if (pool[0].cnt < k) return -1;
    int u = 0, ans = 0;
    for (int i = B - 1; i >= 0; --i) {
      int left = pool[u].nxt[0] ? pool[pool[u].nxt[0]].cnt : 0;
      if (k <= left) u = pool[u].nxt[0];
      else k -= left, u = pool[u].nxt[1], ans |= (1 << i);
    }
    return ans;
  }
  int get_xor_kth(int x, int k) { if (pool[0].cnt < k) return -1;
    int u = 0, ans = 0;
    for (int i = B - 1; i >= 0; --i) {
      int d = (x >> i) & 1;
      int req = pool[u].nxt[d] ? pool[pool[u].nxt[d]].cnt : 0;
      if (k <= req)
        u = pool[u].nxt[d];
      else k -= req, u = pool[u].nxt[d ^ 1],
           ans |= (1 << i);
    }
    return ans;
  }
  int64_t go(int u, int b, int x, int t1, int t2) {
    if (!t1 || t2 == 2) return 0;
    if ((t1 == 2 && !t2) || b < 0) return pool[u].cnt;
    int d = (x >> b) & 1; int64_t res = 0;
    if (d) {
      if (pool[u].nxt[1]) res += go(pool[u].nxt[1], b - 1, x,
               t1 == 1 ? 0 : t1, t2 == 1 ? 0 : t2);
      if (pool[u].nxt[0]) res += go(pool[u].nxt[0], b - 1, x,
               t1 == 1 ? 2 : t1, t2);
    } else {
      if (pool[u].nxt[1]) res += go(pool[u].nxt[1], b - 1, x, t1,
               t2 == 1 ? 2 : t2);
      if (pool[u].nxt[0]) res += go(pool[u].nxt[0], b - 1, x, t1, t2);
    }
    return res;
  }
  int64_t count_special(int x) {
    return go(0, B - 1, x, 1, 1); }
};
\end{minted}
\subsection{Mo's Algorithm}
\texttt{add(pos)} and \texttt{rem(pos)} must be implemented per problem. Snake-order (even/odd block sorting) for $O(N \sqrt{Q})$. Sort queries first. Block size $B = \sqrt{N}$. Alternating sort order within even/odd blocks cuts the total right-pointer moves nearly in half in practice.
\begin{minted}{c++}
// sort queries: blocks of size B, snake-order within blocks
sort(cu.begin(), cu.end(), [&](auto &x, auto &y){
  if (x.l / B != y.l / B)
    return x.l / B < y.l / B;
  return (x.l / B) % 2 ? x.r > y.r : x.r < y.r;
});
int R = -1, L = 0; // [L,R] is current active window
while (R < r) add(++R); // expand right first
while (L > l) add(--L); // expand left
while (R > r) rem(R--); // shrink right
while (L < l) rem(L++); // shrink left
\end{minted}
\subsection{Rollback DSU}
DSU with undo. No path compression --- union by size only. \texttt{time()} returns a snapshot marker; \texttt{rollback(t)} reverts to that state.
\textbf{When to use:} offline dynamic connectivity, DFS with backtracking, problems requiring DSU with rollback.
\begin{minted}{c++}
struct RollbackDSU {
  vector<int> e;
  vector<pair<int,int>> hist;
  RollbackDSU(int n) : e(n, -1) {}
  int find(int x) {
    while (e[x] >= 0) x = e[x];
    return x;
  }
  int size(int x) { return -e[find(x)]; }
  int time() { return hist.size(); }
  void rollback(int t) {
    while ((int)hist.size() > t) {
      auto [x, v] = hist.back(); hist.pop_back();
      e[x] = v;
    }
  }
  bool unite(int a, int b) {
    a = find(a), b = find(b);
    if (a == b) return false;
    if (e[a] > e[b]) swap(a, b);
    hist.push_back({a, e[a]});
    hist.push_back({b, e[b]});
    e[a] += e[b]; e[b] = a;
    return true;
  }
};
\end{minted}
\subsection{Disjoint Sparse Table}
Non-idempotent $O(1)$ range query (sum, xor, AND, OR). Build $O(N \log N)$, query $O(1)$. No updates. Pads to power of 2 internally.
\textbf{When to use:} range sum/xor/AND queries with $O(1)$ time, complement to regular sparse table.
\begin{minted}{c++}
struct DST {
  int n, log;
  vector<vector<int>> t;
  DST(vector<int> a) : n(a.size()) {
    if (n == 0) { log = 0; return; }
    int m = 1;
    while (m < n) m <<= 1;
    a.resize(m, 0);
    log = 32 - __builtin_clz(m);
    t.assign(log, vector<int>(m));
    t[0] = a;
    for (int h = 1; h < log; h++) {
      int half = 1 << (h - 1);
      for (int i = 0; i < m; i += half * 2) {
        int mid = i + half;
        t[h][mid - 1] = a[mid - 1];
        for (int j = mid - 2; j >= i; j--)
          t[h][j] = t[h][j + 1] + a[j];
        t[h][mid] = a[mid];
        for (int j = mid + 1; j < i + half * 2; j++)
          t[h][j] = t[h][j - 1] + a[j];
      }
    }
  }
  int query(int l, int r) {
    if (l == r) return t[0][l];
    int h = 32 - __builtin_clz(l ^ r);
    return t[h][l] + t[h][r];
  }
};
\end{minted}
% For min/max: replace + with min/max, change resize(0) to INF/-INF.
\subsection{Min Queue}
Sliding window minimum in $O(1)$ amortized per op. Maintains deque with increasing values.
Works for max by changing the comparison.
\textbf{When to use:} sliding window min/max, monotonic queue problems, subsequence constraints.
\begin{minted}{c++}
struct MinQueue {
  deque<int> dq; // increasing values
  queue<int> q;  // actual elements
  void push(int x) {
    q.push(x);
    while (!dq.empty() && dq.back() > x) dq.pop_back();
    dq.push_back(x);
  }
  void pop() {
    if (q.front() == dq.front()) dq.pop_front();
    q.pop();
  }
  int get_min() { return dq.front(); }
  bool empty() { return q.empty(); }
};
// For max: change > to < in push, rename to MaxQueue.
\end{minted}
\subsection{Small-to-Large Merging}
Generic technique: when merging two sets, always merge the smaller into the larger.
Total cost $O(N \log N)$ across all merges.
\textbf{When to use:} maintaining per-component data during DSU union or tree DP.
\textbf{Memory tip:} Use \texttt{vector<int>} instead of \texttt{set<int>} when duplicates are not needed --- \texttt{vector} saves $7-8\times$ memory (4 bytes vs $\sim 32$ bytes for \texttt{set} tree nodes). Sort + unique after merge if needed.
\begin{minted}{c++}
// Example: DSU + vector per component (use vector not set)
vector<vector<int>> comp(n); // vector: ~8x less memory than set
vector<int> par(n);
iota(par.begin(), par.end(), 0);
function<int(int)> find = [&](int x) {
  while (par[x] != x) x = par[x] = par[par[x]];
  return x;
};
auto unite = [&](int u, int v) {
  u = find(u), v = find(v);
  if (u == v) return;
  if (comp[u].size() < comp[v].size()) swap(u, v);
  par[v] = u;
  comp[u].insert(comp[u].end(),
    comp[v].begin(), comp[v].end());
  comp[v].clear();
  // sort(comp[u].begin(), comp[u].end());
  // comp[u].erase(unique(...), comp[u].end());
};
\end{minted}
